% Image credits for the photographs used in this volume (Book 1).
% Machine-readable license records live in images/book1/CREDITS.md; keep
% the two files in sync when adding or removing a photograph.
\chapter*{चित्र-श्रेय}

इस पुस्तक के रेखाचित्र मौलिक हैं। छायाचित्र अपने-अपने मुक्त अनुज्ञापत्रों के
अंतर्गत प्रयुक्त हैं; उनके रचयिताओं के प्रति आभार:

\begin{itemize}
  \item \emph{जॉन डाल्टन} (परमाणु और अणु का अध्याय): Charles Turner की
        नक़्क़ाशी, Joseph Allen के चित्र पर आधारित, 1834, Library of Congress,
        Wikimedia Commons के माध्यम से, सार्वजनिक अधिकार-क्षेत्र।
  \item \emph{डाल्टन की तत्व-सारणी, 1808} (परमाणु और अणु का अध्याय):
        John Dalton की पुस्तक \emph{A New System of Chemical Philosophy} का
        एक पृष्ठ, Internet Archive का स्कैन, Wikimedia Commons के माध्यम से,
        सार्वजनिक अधिकार-क्षेत्र।
  \item \emph{बॉक्साइट} (कच्चे माल का अध्याय): Werner Schellmann का
        छायाचित्र, Wikimedia Commons के माध्यम से, CC~BY-SA~2.5।
  \item \emph{ग्रेनाइट} (शुद्ध पदार्थों का अध्याय): Eurico Zimbres का
        छायाचित्र, Wikimedia Commons के माध्यम से, CC~BY-SA~2.0~br।
  \item \emph{निर्जल कॉपर सल्फ़ेट} (पदार्थों की पहचान का अध्याय): W.~Oelen का
        छायाचित्र, Wikimedia Commons के माध्यम से, CC~BY-SA~3.0।
  \item \emph{कॉपर सल्फ़ेट के क्रिस्टल} (पदार्थों की पहचान का अध्याय):
        Stephanb का छायाचित्र, Wikimedia Commons के माध्यम से, CC~BY-SA~3.0।
  \item \emph{आंत्वान-लोरां लावोज़िए और मारी-आन पोल्ज़ लावोज़िए}
        (संतुलित समीकरणों का अध्याय): Jacques-Louis David का तैलचित्र, 1788,
        Metropolitan Museum of Art, Wikimedia Commons के माध्यम से, सार्वजनिक
        अधिकार-क्षेत्र।
  \item \emph{बैकेलाइट का टेलीफ़ोन} (प्लास्टिक का अध्याय): Holger Ellgaard का
        छायाचित्र, Wikimedia Commons के माध्यम से, CC~BY-SA~3.0।
  \item \emph{अर्नेस्ट रदरफ़ोर्ड} (परमाणु के भीतर का अध्याय): छायाचित्र,
        George Grantham Bain Collection, Library of Congress, Wikimedia
        Commons के माध्यम से, सार्वजनिक अधिकार-क्षेत्र।
  \item \emph{हैलाइट} (आयनों का अध्याय): Robert M.~Lavinsky का छायाचित्र,
        Wikimedia Commons के माध्यम से, CC~BY-SA~3.0।
  \item \emph{स्टैच्यू ऑफ़ लिबर्टी} (धातुओं और संक्षारण का अध्याय): Elcobbola का
        छायाचित्र, Wikimedia Commons के माध्यम से, सार्वजनिक अधिकार-क्षेत्र।
  \item \emph{दिमित्री मेंडेलीव} और \emph{उनकी 1869 की सारणी} (आवर्त सारणी का
        अध्याय): Wikimedia Commons के माध्यम से, सार्वजनिक अधिकार-क्षेत्र।
  \item \emph{क्रैब नेबुला} (तत्वों का अध्याय): NASA, ESA, J.~Hester और
        A.~Loll, सार्वजनिक अधिकार-क्षेत्र।
  \item \emph{सोडियम} (इलेक्ट्रॉन कोशों का अध्याय): Dnn87 का छायाचित्र,
        Wikimedia Commons के माध्यम से, CC~BY-SA~3.0।
  \item \emph{एम्प्यूल में क्लोरीन} (इलेक्ट्रॉन कोशों का अध्याय): W.~Oelen का
        छायाचित्र, Wikimedia Commons के माध्यम से, CC~BY-SA~3.0।
  \item \emph{आमेदेओ आवोगाद्रो} (मोल का अध्याय): C.~Sentier के एक रेखाचित्र पर
        आधारित लिथोग्राफ़, 1856, Edgar Fahs Smith collection, Wikimedia
        Commons के माध्यम से, सार्वजनिक अधिकार-क्षेत्र।
  \item \emph{हिम-क्रिस्टल} (ध्रुवता का अध्याय): Wilson Bentley के छायाचित्र,
        1902, Wikimedia Commons के माध्यम से, सार्वजनिक अधिकार-क्षेत्र।
  \item \emph{लुई पाश्चर} (त्रिविम रसायन का अध्याय): Paul Nadar का छायाचित्र,
        Wikimedia Commons के माध्यम से, सार्वजनिक अधिकार-क्षेत्र।
  \item \emph{अमोनियम टार्टरेट के क्रिस्टल} (त्रिविम रसायन का अध्याय):
        Marie-Lan Taÿ Pamart का छायाचित्र, Wikimedia Commons के माध्यम से,
        CC~BY~4.0।
  \item \emph{स्वांते आरेनियस} (अम्लों और क्षारकों का अध्याय): फ़ोटोग्रेव्योर,
        1909, Wikimedia Commons के माध्यम से, सार्वजनिक अधिकार-क्षेत्र।
  \item \emph{वोल्टा का पुंज} (विद्युत-रासायनिक सेलों का अध्याय): GuidoB का छायाचित्र,
        Wikimedia Commons के माध्यम से, CC~BY-SA~3.0।
  \item \emph{अपनी प्रयोगशाला में वॉलेस कैरोथर्स} (बहुलकों का अध्याय):
        छायाकार अज्ञात, Wikimedia Commons के माध्यम से, सार्वजनिक
        अधिकार-क्षेत्र।
\end{itemize}

हर छायाचित्र के पूरे स्रोत-लिंक और अनुज्ञापत्रों के URL पुस्तक के भंडार में
\texttt{images/book1/CREDITS.md} में सूचीबद्ध हैं। ख़तरे के चित्रसंकेत रसायनों के
वर्गीकरण और लेबलिंग की विश्वव्यापी सामंजस्यपूर्ण प्रणाली (GHS) के हैं, जैसे वे
\TeX\ Live के \texttt{ghsystem} पैकेज के लिए बनाए गए हैं।

\bigskip
छायाचित्रों और मौलिक रेखाचित्रों के साथ-साथ कुछ चित्र AI से बनाए गए हैं: वे
पुस्तक के संपादकों द्वारा लिखे गए निर्देशों से OpenAI के चित्र-निर्माण द्वारा बनाए
गए, और शामिल करने से पहले रासायनिक शुद्धता के लिए जाँचे गए। हर बनाए गए चित्र का
पूरा निर्देश भंडार में \texttt{images/book1/ai/PROMPTS.md} में दर्ज है।
\clearpage
